\documentclass[dvipdfmx,11pt,a4paper]{article}
\usepackage{amsmath}
\usepackage{amssymb}
\usepackage{amsthm}
\usepackage{amsfonts}
\usepackage{bm}
\usepackage[utf8]{inputenc}
\usepackage[dvipdfmx]{geometry}
\geometry{left=25mm,right=25mm,top=25mm,bottom=25mm}

\title{二重時空のねじれ不整合の $G(3,1,1)$ 拡張：\\
10次元生成子構造と面に基づくブラケット積}

\author{https://github.com/hypernumbernet}
\date{\today}

\theoremstyle{definition}
\newtheorem{definition}{定義}
\newtheorem{conjecture}{予想}
\theoremstyle{plain}
\newtheorem{theorem}{定理}
\theoremstyle{remark}
\newtheorem{remark}{注意}

\begin{document}

\maketitle

\begin{abstract}
二重時空構造を、$\mathrm{Cl}(3,1)$ の双四元数基底に零ベクトル $e_{4}$ を付加することにより射影幾何代数 $G(3,1,1)$ に埋め込む。得られるねじれ不整合は、3つの双曲型生成子（$iI, iJ, iK$）、3つの巡回型生成子（$I, J, K$）、および $e_{4}$ と通常時空の4つのベクトル生成子 $\{j, kI, kJ, kK\}$ から構成される4つの零生成子 $N_{\mu} := e_{4}\wedge e_{\mu}$（$\mu=0,1,2,3$）からなる10次元バイベクトル空間によって生成される。この空間はポアンカレ代数 $\mathfrak{iso}(3,1)\cong\mathfrak{so}(3,1)\ltimes\mathbb{R}^{3,1}$ の構造的候補である：6つの双曲・巡回バイベクトルはローレンツコピー $\mathfrak{so}(3,1)$ を張り、4つの零バイベクトルは幾何積の下で可換な並進イデアルを形成する。

ねじれプラス並進パラメーター上の二次形式 $J^{(5)}$ は、6つのねじれ係数上のカルタン--キリング二次と並進係数上のミンコフスキー計量を組み合わせたものであり、パラメーター空間上のスカラー診断として用いる。3ブレード上の面に基づくブラケット積と双対＝法線対応は開いた幾何プログラムとして残す（付録Bは代表計算を与える）。各零生成子 $N_{\mu}$ は平方が零で反転奇（$\widetilde{N_{\mu}}=-N_{\mu}$）であるバイベクトルであるため、積 $M:=RT$（$R=\exp(\Omega_{\mathrm{torsion}})$、$T=1+\Omega_{\mathrm{trans}}$）で定義されるモーターはユニタリーであり $M\widetilde{M}=1$ となる。$[\Omega_{\mathrm{torsion}},\Omega_{\mathrm{trans}}]\neq 0$ のとき、この積と $\exp(\Omega_{\mathrm{biv}}^{(5)})$ との関係は後続の課題とする。
\end{abstract}
 
\section{はじめに}
\label{sec:introduction}

二重時空枠組みは、各質量粒子が内在的にペアをなす2つのミンコフスキー構造の間のねじれ不整合として重力と慣性を記述する。並進自由度を幾何的に自然な形で取り込むため、生成ベクトル空間が5次元でクリフォード代数全体が $2^{5}=32$ 次元である射影幾何代数 $G(3,1,1)$~\cite{Gunn2017,Lengyel2024,DeKeninckDorst2019,GunnDeKeninck2020} に構造を持ち上げる。平方が零の退化ベクトルを用いて並進をバイベクトルとして符号化する手法は、共形／射影クリフォード代数における Selig の剛体運動学の扱い~\cite{Selig2005} に遡り、現在ではユークリッドおよびミンコフスキー符号の両方において標準的な PGA 慣習となっている~\cite{Lengyel2024,GunnDeKeninck2020}。付加的な零ベクトル $e_{4}$ は4つの $\mathrm{Cl}(3,1)$ ベクトル生成子 $\{j, kI, kJ, kK\}$~\cite{Hestenes2015,DoranLasenby2003} と反可換であり $e_{4}^{2}=0$ を満たす。さらに双四元数擬スカラー $i=jkI\,kJ\,kK$~\cite{ConwaySmith2003} と可換であるため、双対写像 $X\mapsto Xi$ は $\mathrm{Cl}(3,1)$ 部分代数上では親論文と同様に作用しつつ、新しい零方向は不変のままである。

不整合バイベクトルは3つの異なる類に分かれる。3つの双曲型生成子 $iI, iJ, iK$ は平方が $+1$ で $\mathfrak{so}(3,1)$ のブーストセクターを張る。3つの巡回型生成子 $I, J, K$ は平方が $-1$ で、同じローレンツコピーの回転セクターを張る（DST 双対は符号まで両セクターを入れ替える）。4つの零生成子 $N_{\mu} := e_{4}\wedge e_{\mu}$（$\mu=0,1,2,3$）は $e_{4}$ と4つの $\mathrm{Cl}(3,1)$ ベクトル生成子から構成され、平方は零で可換な並進イデアルを形成する。これにより10次元バイベクトル空間が得られ、ポアンカレ代数 $\mathfrak{iso}(3,1)\cong\mathfrak{so}(3,1)\ltimes\mathbb{R}^{3,1}$~\cite{Weinberg1995,Tung1985} の次元と一致する。一致は構造的候補として扱い、完全なリー括弧同型は本稿の範囲を超える。したがって並進は補助ゲージ場なしに零セクターから直接現れ、得られる拡張不整合バイベクトルは、単一の代数対象の中で元のねじれ自由度と並進自由度を統一する。

第二の主題は $G(3,1,1)$ の面幾何である。3ブレード上の面に基づくブラケット積は、双対時空を法線束として読む幾何解釈を示唆する；付録Bには代表計算を示し、一般の法線写像定理は開のままである。

各零バイベクトルは平方が零である \emph{かつ} 任意の $\mu,\nu$ に対して $N_{\mu}N_{\nu}=0$ を満たすため、零部分の指数関数は1次の有限多項式である。すべてのバイベクトルの反転奇性 $\widetilde{N_{\mu}}=-N_{\mu}$ と合わせて、これは積 $M:=RT$ で定義されるモーターのユニタリー性 $M\widetilde{M}=1$ を保証する。$G(3,1,1)$ の退化計量全体に関するサンドイッチ、および因子が可換でないときの $M$ と $\exp(\Omega_{\mathrm{biv}}^{(5)})$ との同一視はここでは扱わない。

本稿の構成は次のとおりである。第2節では二重時空の $G(3,1,1)$ への埋め込みと10生成子の分類を定義する。第3節では拡張ねじれ不整合バイベクトルを構成し、零指数の多項式形を証明し、パラメーター空間形式 $J^{(5)}$ を導入する。第4節では面に基づくブラケットプログラムの概要を述べる。第5節では積モーター $M:=RT$ のユニタリー性を確立する。第6節では固定チャート上の PGA--TEGR 対応を展開し、シュワルツシルト対角 ansatz と $J^{(5)}\leftrightarrow T$ 同一視の位置づけを述べる。第7節で結論を述べる。

\section{10次元生成子構造}

\subsection{ベクトル基底と $e_{4}$ の付加}

親理論では4つの $\mathrm{Cl}(3,1)$ ベクトル生成子を双四元数の元として
\[
e_{0}=j,\qquad e_{1}=kI,\qquad e_{2}=kJ,\qquad e_{3}=kK,
\]
と同一視し、ミンコフスキー符号 $(-,+,+,+)$ を満たす：
\[
e_{0}^{2}=-1,\qquad e_{1}^{2}=e_{2}^{2}=e_{3}^{2}=+1,\qquad \{e_{\mu},e_{\nu}\}=0\ (\mu\neq\nu).
\]
二重時空枠組みのねじれ的性質を保ちつつ並進自由度を取り込むため、この構造を単一の新しいベクトル生成子 $e_{4}$ を付加することで $G(3,1,1)$ に持ち上げ、
\begin{equation}
e_{4}^{2}=0,\qquad \{e_{4},e_{\mu}\}=0\quad (\mu=0,1,2,3).
\label{eq:e4-defs}
\end{equation}
で特徴づける。完全なクリフォード代数 $G(3,1,1)$ の次元は $2^{5}=32$ である；元の16次元双四元数代数は $\{e_{0},e_{1},e_{2},e_{3}\}$ で生成される部分代数として回復される。

\subsection{双四元数双対写像との両立}

親理論の双四元数擬スカラーは
\[
i = e_{0}e_{1}e_{2}e_{3} = j\,(kI)(kJ)(kK),
\]
であり、二重時空の双対写像は $X\mapsto Xi$ として作用する。$\{e_{4},e_{\mu}\}=0$ を繰り返し用いると
\[
e_{4}\,i \;=\; e_{4}\,e_{0}e_{1}e_{2}e_{3} \;=\; (-1)^{4}\,e_{0}e_{1}e_{2}e_{3}\,e_{4} \;=\; i\,e_{4},
\]
となり $e_{4}$ は $i$ と可換である。したがって双対写像 $X\mapsto Xi$ は $\mathrm{Cl}(3,1)$ 部分代数上では親論文と全く同様に作用し、$e_{4}$ は独立な零方向として不変のままである。これが双対を「$e_{4}$ の零性を保ちつつ」$G(3,1,1)$ に拡張するという意味での正確な内容である。

\subsection{10バイベクトル生成子の分類}

$G(3,1,1)$ のバイベクトル空間の次元は $\binom{5}{2}=10$ である。6つの $\mathrm{Cl}(3,1)$ バイベクトルと $e_{4}$ を含む4つのバイベクトルを合わせ、拡張ねじれ不整合の生成子はその平方の符号に応じて3つの異なる類に分かれる：

\begin{itemize}
  \item \textbf{双曲型生成子}（平方 $+1$）：
  \[
  iI = e_{0}e_{1},\qquad iJ = e_{0}e_{2},\qquad iK = e_{0}e_{3}.
  \]
  これら3元は $\mathfrak{so}(3,1)$ のブーストセクターを張る。

  \item \textbf{巡回型生成子}（平方 $-1$）：
  \[
  I = e_{3}e_{2},\qquad J = e_{1}e_{3},\qquad K = e_{2}e_{1}.
  \]
  これら3元は同じローレンツコピー $\mathfrak{so}(3,1)$ の回転セクターを張る。

  \item \textbf{零型生成子}（平方 $0$）：$e_{4}$ と4つの $\mathrm{Cl}(3,1)$ ベクトル生成子から構成される4つのバイベクトル、
  \[
  N_{0} = e_{4}\wedge e_{0} = e_{4}\,j,\qquad
  N_{1} = e_{4}\wedge e_{1} = e_{4}\,(kI),
  \]
  \[
  N_{2} = e_{4}\wedge e_{2} = e_{4}\,(kJ),\qquad
  N_{3} = e_{4}\wedge e_{3} = e_{4}\,(kK).
  \]
  \eqref{eq:e4-defs} より、任意の $\mu,\nu\in\{0,1,2,3\}$ に対して
  \begin{equation}
  N_{\mu}N_{\nu} \;=\; (e_{4}e_{\mu})(e_{4}e_{\nu}) \;=\; -\,e_{4}^{2}\,e_{\mu}e_{\nu} \;=\; 0,
  \label{eq:null-strong-vanishing}
  \end{equation}
  となり、各 $N_{\mu}$ の平方は零（$\mu=\nu$ の場合）であり、\emph{零生成子の任意の組の積}も消える。4つの $N_{\mu}$ は10次元代数の可換な並進イデアルを張る。
\end{itemize}

これらの類全体で10次元バイベクトル空間が与えられる：
\[
\dim \;=\; 3 \text{（双曲）} \;+\; 3 \text{（巡回）} \;+\; 4 \text{（零）} \;=\; 10,
\]
ポアンカレ代数 $\mathfrak{iso}(3,1)$ の次元と一致する。ローレンツ部分 $\mathfrak{so}(3,1)$ は6つの双曲・巡回バイベクトルによって生成され、4つの零バイベクトル $N_{\mu}$ は幾何積の下で並進生成子 $P_{\mu}$ を可換イデアルとして実現する。したがってこの構成は $\mathfrak{so}(3,1)\ltimes\mathbb{R}^{3,1}$ の次元・構造上の候補であり、完全なリー同型定理はここでは展開しない。

\subsection{零セクターと並進}

零セクターは時空並進を直接符号化する。実際、任意の並進パラメーター $a^{\mu}$（$\mu=0,1,2,3$）に対し、和 $T(a):=\tfrac{1}{2}a^{\mu}N_{\mu}$（アインシュタイン和を暗黙のうちに）は零セクターだけから構成されるバイベクトルである。\eqref{eq:null-strong-vanishing} より $T(a)^{2}=0$ なので指数
\[
\exp\bigl(T(a)\bigr) \;=\; 1 + \tfrac{1}{2}\,a^{\mu}N_{\mu}
\]
は1次で打ち切られる。したがって並進作用は純粋なシフトとして現れ、追加の曲率や高次項は入らない；特に光様並進は4次元パラメーター空間内の錐 $\eta_{\mu\nu}a^{\mu}a^{\nu}=0$ に対応する。

\subsection{零生成子上の双対性}

双対写像 $X\mapsto Xi$ は $\mathrm{Cl}(3,1)$ 部分代数上で親論文と同様に作用するため、ブースト生成子 $\{iI,iJ,iK\}$ と回転生成子 $\{I,J,K\}$ を符号まで交換する。さらに $e_{4}$ は $i$ と可換なので
\[
N_{\mu}\,i \;=\; (e_{4}e_{\mu})\,i \;=\; e_{4}\,(e_{\mu}\,i).
\]
ここで $e_{\mu}i$ は Cl(3,1) 部分代数で次数3であり、したがって $N_{\mu}i$ は $G(3,1,1)$ で次数4である。ゆえに双対は各零生成子をバイベクトル次数の外へ写す：4次元の零バイベクトルイデアルは双対閉ではなく、一方6次元ねじれセクターは次数2の内部で双対の下に閉じる。

要約すると、$G(3,1,1)$ への埋め込みは拡張不整合をポアンカレ型 $\mathfrak{so}(3,1)\ltimes\mathbb{R}^{3,1}$ の10次元バイベクトル空間（構造的候補として）に整理し、並進は $e_{4}$ と通常時空ベクトル生成子から構成される4次元零セクターから自然に現れる。

\section{拡張ねじれ不整合の構成}

前節で導入した10次元生成子空間上の拡張ねじれ不整合バイベクトルを構成する。$\Omega_{\text{biv}}^{(5)}$ を拡張モーターの無限小生成子とする。3類の生成子に従って分解する：

\begin{equation}
\Omega_{\text{biv}}^{(5)}
= \underbrace{\sum_{a=1}^{3} \Bigl( \frac{\alpha_{a}}{2}\, B_{a}^{+} + \frac{\beta_{a}}{2}\, B_{a}^{-} \Bigr)}_{\Omega_{\text{torsion}}\ :\ \text{双曲+巡回（6次元）}}
\;+\;
\underbrace{\sum_{\mu=0}^{3} \frac{\lambda^{\mu}}{2}\, N_{\mu}}_{\Omega_{\text{trans}}\ :\ \text{零セクター（4次元）}},
\label{eq:Omega-decomp}
\end{equation}

ここで $B_{a}^{+}$ は双曲型生成子（$iI, iJ, iK$）、$B_{a}^{-}$ は巡回型生成子（$I, J, K$）、$N_{\mu}=e_{4}\wedge e_{\mu}$（$\mu=0,1,2,3$）は第2節で構成した4つの零バイベクトルである。

係数 $\alpha_{a}$ と $\beta_{a}$ は通常セクターと双対セクター間の元のねじれ不整合をパラメーター化し、4係数 $\lambda^{\mu}$ はローレンツ4ベクトルとしての並進不整合を符号化する。光様並進は4次元パラメーター空間内の錐 $\eta_{\mu\nu}\lambda^{\mu}\lambda^{\nu}=0$ に対応する。

\subsection{零指数の厳密な打ち切り}

強い消去性 \eqref{eq:null-strong-vanishing} により $N_{\mu}N_{\nu}=0$ が \emph{すべて} の組 $(\mu,\nu)$---単に $\mu=\nu$ だけではなく---に対して成り立つので
\begin{equation}
\left(\sum_{\mu=0}^{3}\frac{\lambda^{\mu}}{2}N_{\mu}\right)^{\!\!2}
\;=\;\frac{1}{4}\sum_{\mu,\nu=0}^{3}\lambda^{\mu}\lambda^{\nu}\,N_{\mu}N_{\nu}
\;=\;0,
\end{equation}
となり、零部分の指数は1次で打ち切れる：
\begin{equation}
\exp\!\Bigl( \Omega_{\text{trans}} \Bigr)
\;=\;\exp\!\Bigl(\sum_{\mu=0}^{3}\frac{\lambda^{\mu}}{2}N_{\mu}\Bigr)
\;=\;1+\sum_{\mu=0}^{3}\frac{\lambda^{\mu}}{2}N_{\mu}.
\label{eq:exp-truncation}
\end{equation}
\eqref{eq:exp-truncation} の成立には \emph{すべて} の混合積 $N_{\mu}N_{\nu}$ が消えることが必要であり、対角の平方 $N_{\mu}^{2}=0$ だけでは足りない；強い消去 \eqref{eq:null-strong-vanishing} がまさにこれを与える。

\subsection{モーターの因子分解}

各 $N_{\mu}$ が $N_{\mu}N_{\nu}=0$ を満たすため、4つの零生成子は可換である。一方交差交換子 $[B_{a}^{\pm},N_{\mu}]$ は一般に非零であり、一般には $\Omega_{\text{torsion}}$ と $\Omega_{\text{trans}}$ は可換でない。したがって順序付きの積で定義されるモーター
\begin{equation}
M \;:=\; R\,T,
\qquad
R := \exp\!\bigl(\Omega_{\text{torsion}}\bigr),\qquad
T := \exp\!\bigl(\Omega_{\text{trans}}\bigr)=1+\Omega_{\text{trans}},
\end{equation}
を用い、これに対して第5節で $M\widetilde{M}=1$ を証明する。因子が可換でないとき、$\exp(\Omega_{\text{biv}}^{(5)})$ は一般に $RT$ と異なる；その同一視（および共役による再パラメーター化 $T=\exp(\Omega'_{\text{trans}})$）は先送りする。

\subsection{双対性の再考}

双曲成分と巡回成分は親論文と同様に双対の下で互いに変形する。第2節で記録したとおり、双対は各零生成子 $N_{\mu}$ を次数4の元へ送るので、並進パラメーターは双対化されたねじれセクターとは別に追跡すべきである。

\subsection{10次元空間上の二次不変量}

ポアンカレ代数 $\mathfrak{iso}(3,1)$ 上のカルタン--キリング形式は \emph{退化} していることは衆知である~\cite{FultonHarris1991,Humphreys1972}：並進生成子 $P_{\mu}$ に対して
\[
B_{\text{Killing}}(P_{\mu},P_{\nu})\;=\;0.
\]
したがってキリング形式単独では並進不整合を検出する二次不変量を与えられない。そこでローレンツセクターのキリング形式を並進セクター上のミンコフスキー計量で拡張し、10次元生成子空間上の非退化双線形形式を得る。明示的に
\begin{equation}
\boxed{\;
J^{(5)}
\;:=\;
\tfrac{1}{16}\,B_{\text{Killing}}\!\bigl(\Omega_{\text{torsion}},\Omega_{\text{torsion}}\bigr)
\;+\;
\tfrac{1}{2}\,\eta_{\mu\nu}\,\lambda^{\mu}\lambda^{\nu},
\;}
\label{eq:J5-invariant}
\end{equation}
と置く。ここで $B_{\text{Killing}}$ は親論文で用いた6次元ねじれ（ローレンツ）係数上のカルタン--キリング型二次形式であり、$\eta_{\mu\nu}=\mathrm{diag}(-1,+1,+1,+1)$ は並進セクター上のミンコフスキー計量である。第1項は親理論のねじれスカラー $J$ を再現し、第2項は並進が零であるときに消える光錐 $\eta_{\mu\nu}\lambda^{\mu}\lambda^{\nu}=0$ 上で正確に零となる、ローレンツ不変な並進不整合の二次測度を与える。2つの寄与は $\mathfrak{iso}(3,1)$ 上のキリング形式が半単純と並進の間で消えるため結合しない。

\begin{remark}[\(J^{(5)}\) の有界性]
\label{rem:J5-bounded}
並進寄与 $\tfrac{1}{2}\eta_{\mu\nu}\lambda^{\mu}\lambda^{\nu}$ は $\lambda^{\mu}$ を非制約にしたとき有界ではない：消えるのは零条件 $\eta_{\mu\nu}\lambda^{\mu}\lambda^{\nu}=0$ のみである。離散ディオファントス同伴論文の有界性定理は、許容構成上のねじれ部分 $J$（正規化スカラー $J_{\rm norm}=\frac{8}{3\pi^2}J$ と等価）に適用され、$J^{(5)}$ 全体へ自動的には拡張されない。
\end{remark}

このように、拡張ねじれ不整合バイベクトル $\Omega_{\text{biv}}^{(5)}$ は元の6次元ねじれ自由度と4次元並進セクターを単一の10次元代数対象の中に統一し、自然な非退化二次不変量 \eqref{eq:J5-invariant} が両セクターに結合する。

\section{面に基づくブラケット積}
\label{sec:brackets-open}

\subsection{ブラケット積とその次数内容}

$G(3,1,1)$ では面は3ブレードで表される。2つの3ブレード $A,B$ に対し、\emph{面に基づくブラケット積}を幾何積の反対称部分と対称部分として定義する：
\begin{equation}
A \mathbin{\underline{\vee}} B \;:=\; \tfrac{1}{2}(AB - BA),
\qquad
A \mathbin{\overline{\vee}} B \;:=\; \tfrac{1}{2}(AB + BA),
\label{eq:bracket-def}
\end{equation}
右辺の積は代数の幾何積である。\footnote{記号 $\underline{\vee}$ と $\overline{\vee}$ は本稿では \eqref{eq:bracket-def} で定義されるブラケット積に用いる；標準的な射影幾何代数~\cite{Lengyel2024,DorstFontijneMann2007} の「反内積」「反外積」と混同してはならない。}

5ベクトル代数 $G(3,1,1)$ 上の次数3の斉次マルチベクトル2つに対し、幾何積の次数内容は
\[
A_{3}B_{3} \;=\; \langle A_{3}B_{3}\rangle_{0} \;+\; \langle A_{3}B_{3}\rangle_{2} \;+\; \langle A_{3}B_{3}\rangle_{4},
\]
である（5次元ベクトル空間には次数6成分が存在しない）。付録Bは対称性パターン
\[
\langle A_{3}B_{3}\rangle_{0} = \langle B_{3}A_{3}\rangle_{0},
\qquad
\langle A_{3}B_{3}\rangle_{2} = -\,\langle B_{3}A_{3}\rangle_{2},
\qquad
\langle A_{3}B_{3}\rangle_{4} = \langle B_{3}A_{3}\rangle_{4}
\]
を代表例で検証する；任意の3ブレードに対する定理は開のままである。

\subsection{双対＝法線対応}

$P$ を通常時空ベクトル生成子から構成した3ブレード、$D$ を双対時空セクターの元とする。自然な予想は、$D$ が $P$ の法線方向に双対であるとき、
\begin{equation}
P \mathbin{\underline{\vee}} D \;=\; \kappa \, (e_{4}\wedge n),
\label{eq:normal-correspondence}
\end{equation}
が法線ベクトル $n$ と係数 $\kappa$ に対して成り立つというものである。付録Bは具体例を記録する；一般の主張は開のままである。

\subsection{反交換子による並進不整合}

零寄与を含む $P\mathbin{\overline{\vee}}D$ の展開は、$e_{4}\wedge n\wedge e_{\mu}$ 形の次数4担体を示唆する。これらのブラケットを $\Omega_{\text{biv}}^{(5)}$ に結びつける完全な計算は将来の課題である。

\section{モーター $M:=RT$ のユニタリー性}

\subsection{バイベクトル生成子の反転奇性}

$G(3,1,1)$ 上の反転操作 $\,\tilde{}\,$ は次数2の元をその負に送る。実際 $\mu\neq\nu$ に対して $(e_{\mu}e_{\nu})^{\sim} = e_{\nu}e_{\mu} = -e_{\mu}e_{\nu}$ である。特に第2節で導入した10バイベクトル生成子はすべて反転奇である：
\begin{equation}
\widetilde{B_{a}^{+}} = -B_{a}^{+},\qquad
\widetilde{B_{a}^{-}} = -B_{a}^{-},\qquad
\widetilde{N_{\mu}}   = -N_{\mu}.
\label{eq:rev-odd}
\end{equation}
$N_{\mu}$ の反転奇性は $N_{\mu}$ がバイベクトル $e_{4}\wedge e_{\mu}$ であることに本質的に依存する。次数1の零ベクトルを並進生成子に選んでいたら反転偶となり、以下のユニタリー論証は成立しなかったであろう。\footnote{射影幾何代数において並進生成子に零\emph{ベクトル}ではなく零\emph{バイベクトル}を採用する幾何的理由である。ユークリッドおよびミンコフスキー符号の標準PGA文献~\cite{Gunn2017,Lengyel2024,GunnDeKeninck2020} でも同じ慣習が用いられる。}

\subsection{積モーター $M:=RT$ に対するユニタリー性の明示計算}

積モーター $M:=R\,T$ を考える。ねじれローターに対して、反転奇性 \eqref{eq:rev-odd} から標準的なユニタリー性
\[
R \, \widetilde{R} \;=\; \exp(\Omega_{\text{torsion}})\,\exp(\widetilde{\Omega_{\text{torsion}}})
\;=\; \exp(\Omega_{\text{torsion}})\,\exp(-\Omega_{\text{torsion}}) \;=\; 1,
\]
が従い、これは6次元ねじれセクターに制限した親理論の $R\widetilde{R}=1$ である。零翻訳子には \eqref{eq:exp-truncation} と \eqref{eq:rev-odd} を用いる：
\begin{align}
T \, \widetilde{T}
&= \left(1+\sum_{\mu}\frac{\lambda^{\mu}}{2}N_{\mu}\right)\!\left(1+\sum_{\nu}\frac{\lambda^{\nu}}{2}\widetilde{N_{\nu}}\right) \notag\\
&= \left(1+\sum_{\mu}\frac{\lambda^{\mu}}{2}N_{\mu}\right)\!\left(1-\sum_{\nu}\frac{\lambda^{\nu}}{2}N_{\nu}\right) \notag\\
&= 1 \;-\; \underbrace{\left(\sum_{\mu}\frac{\lambda^{\mu}}{2}N_{\mu}\right)^{\!\!2}}_{=\,0\ \text{より}\ \eqref{eq:null-strong-vanishing}}
\;=\; 1.
\label{eq:null-unitarity}
\end{align}
反転は反準同型（$\widetilde{R\,T}=\widetilde{T}\,\widetilde{R}$）なので
\begin{equation}
M\,\widetilde{M}
\;=\; R\,T\,\widetilde{T}\,\widetilde{R}
\;=\; R\,\widetilde{R}
\;=\; 1.
\end{equation}
これにより順序付き積 $M=RT$ のユニタリー性が確立される。$\Omega_{\text{torsion}}$ と $\Omega_{\text{trans}}$ が可換でないとき、$M$ は $\exp(\Omega_{\text{biv}}^{(5)})$ と一致する必要はない；$G(3,1,1)$ の退化計量全体に関するサンドイッチも同様に扱わない。

\subsection{統一的代数記述}

この積モーターの枠組みでは、光様並進（$\eta_{\mu\nu}\lambda^{\mu}\lambda^{\nu}=0$）とねじれ不整合は $J^{(5)}$ で測られる同じ10次元パラメーター空間に位置する。双対が零バイベクトルイデアルを閉じないため（第2節）、並進パラメーターは双対化されたねじれデータとは別に追跡する。

\section{PGA--TEGR 対応}
\label{sec:pga-tegr}
\label{sec:schwarzschild-pga}

\subsection{範囲}

親理論の二重時空はテレパラレル重力の実現を目指す：幾何は非零曲率テンソルではなくねじれによって担われる。射影拡張 $G(3,1,1)$ では、そのプログラムを固定座標チャート上で展開する。終始ポアンカレ候補
$\mathfrak{so}(3,1)\ltimes\mathbb{R}^{3,1}$ を用いる（12次元ラベル $\mathfrak{so}(3,1)\oplus\mathfrak{so}(3,1)$ は用いない）。本節ではモーター誘導テトラッドを構成し、ヴァイツェンベックねじれを記録し、外部シュワルツシルト対角 ansatz に特殊化し、パラメーター空間形式 $J$、$J^{(5)}$ とテレパラレルスカラー $T$ の同一視の位置づけを述べる。一般多様体上での TEGR とアインシュタイン--ヒルベルト作用の完全な変分的等価は TEGR 文献~\cite{maluf2013} から取り、ここでは再導出しない。

\begin{center}
\footnotesize
\setlength{\tabcolsep}{4pt}
\begin{tabular}{@{}p{0.42\textwidth}p{0.26\textwidth}l@{}}
\hline
主張 & 状態 & 形式化 \\
\hline
モーター $M:=RT$ のユニタリー性 & 証明済（第5節） & Lean \texttt{Motor} \\
対角シュワルツシルトテトラッドからの誘導計量 & 証明済（チャート） & Lean \texttt{Gravity.Schwarzschild} \\
外部チャート上のヴァイツェンベック $T^{a}{}_{\mu\nu}$ とスカラー $T$ & 証明済（チャート） & Lean \texttt{Gravity.Weitzenbock} \\
動径ブースト尺度因子がシュワルツシルト赤方偏移と一致 & 証明済 & Lean \texttt{Gravity.Sandwich} \\
静的対角ゲージで $J=\tfrac12 T+\nabla\cdot B$ & 証明済（チャート切片） & Lean \texttt{Gravity.Identification} \\
任意モーターに対する一般の $J^{(5)}\leftrightarrow T$ & 予想 & 下記 \\
TEGR $\Leftrightarrow$ アインシュタイン--ヒルベルト（変分） & 外部入力 & \cite{maluf2013} \\
\hline
\end{tabular}
\end{center}

\subsection{モーター誘導テトラッド}

\begin{definition}[基準共枠とサンドイッチ]
$G(3,1,1)$ 内部の Cl(3,1) ベクトル生成子 $\{\mathring{e}^{a}\}_{a=0}^{3}$ を固定する。ユニタリーモーター $M$（第5節）に対し、サンドイッチされたベクトルを
\begin{equation}
e^{a} \;:=\; \bigl\langle M\,\mathring{e}^{a}\,\widetilde{M}\bigr\rangle_{1}.
\label{eq:motor-tetrad}
\end{equation}
と定義する。座標チャート $(x^{\mu})$ 上では成分場 $e^{a}{}_{\mu}(x)$ と誘導計量
\begin{equation}
g_{\mu\nu} \;=\; \eta_{ab}\, e^{a}{}_{\mu}\, e^{b}{}_{\nu},
\label{eq:induced-metric}
\end{equation}
を書き、$\eta=\mathrm{diag}(-1,+1,+1,+1)$ とする。
\end{definition}

$M=R=\exp(\Omega_{\mathrm{torsion}})$ が純ローレンツローターのとき、サンドイッチは基準ベクトル上の局所ローレンツ作用を実現する。零並進 $T=1+\Omega_{\mathrm{trans}}$ は完全モーター $M=RT$ に入るが、それ自体では古典的シュワルツシルト対角ゲージを与えず、$J^{(5)}$ で測られる並進不整合診断に利用可能のままである。

\subsection{固定チャート上のヴァイツェンベックねじれ}

\begin{definition}[チャート・ヴァイツェンベックねじれ]
固定チャート上、共枠 $e^{a}{}_{\mu}$ のヴァイツェンベックねじれは
\begin{equation}
T^{a}{}_{\mu\nu} \;=\; \partial_{\mu} e^{a}{}_{\nu} - \partial_{\nu} e^{a}{}_{\mu}.
\label{eq:weitzenbock-torsion}
\end{equation}
である。逆テトラッド $e_{a}{}^{\mu}$ と $g_{\mu\nu}$ で下げた時空添字を用い、テレパラレルねじれスカラーは
\begin{equation}
T \;=\; \tfrac14 T_{\rho\mu\nu}T^{\rho\mu\nu} + \tfrac12 T_{\rho\mu\nu}T^{\nu\mu\rho} - T_{\rho}T^{\rho},
\label{eq:teleparallel-T}
\end{equation}
である。ここで $T_{\rho}=T^{\sigma}{}_{\rho\sigma}$。
\end{definition}

この段階では計量適合なレヴィ--チビタ接続の曲率テンソルは不要である：テレパラレル記述は幾何を $T^{a}{}_{\mu\nu}$ に格納する。

\subsection{シュワルツシルト ansatz}

\begin{definition}[外部シュワルツシルト対角テトラッド]
$r_{s}=2GM/c^{2}>0$ とし、外部領域 $r>r_{s}$ で働く。対角共枠（単位 $c=1$）は
\begin{align}
e^{0} &= \sqrt{1-\frac{r_{s}}{r}}\,dt,\nonumber\\
e^{1} &= \bigl(1-\frac{r_{s}}{r}\bigr)^{-1/2}\,dr,\nonumber\\
e^{2} &= r\,d\theta,\nonumber\\
e^{3} &= r\sin\theta\,d\phi.
\label{eq:schwarzschild-tetrad}
\end{align}
である。
\end{definition}

\begin{definition}[動径ブーストラピディティ]
$A(r):=1-r_{s}/r$ と置き、
\begin{equation}
\varphi(r) \;:=\; -\log\sqrt{A(r)} \;=\; -\tfrac12\log A(r)
\qquad(r>r_{s}).
\label{eq:radial-rapidity}
\end{equation}
とする。このとき $e^{-\varphi}=\sqrt{A}$、$e^{\varphi}=A^{-1/2}$ であり、\eqref{eq:schwarzschild-tetrad} の $(t,r)$ 対角係数は $B^{+}_{1}=e_{0}\wedge e_{1}$ で生成される純双曲ローターのブースト尺度因子である。角脚 $e^{2},e^{3}$ は球対称基準共枠から取る。対応するねじれパラメーターは純ブースト $\Omega_{\mathrm{torsion}}=(\varphi/2)B^{+}_{1}$ であり、対角ゲージでは零パラメーター $\lambda^{\mu}=0$ である。
\end{definition}

\subsection{$J$ および $J^{(5)}$ と $T$ の同一視}

\begin{theorem}[チャート水準の主張]
\eqref{eq:schwarzschild-tetrad} をもつ外部チャート上で：
\begin{enumerate}
\item 誘導計量~\eqref{eq:induced-metric} はシュワルツシルト計量である；
\item ヴァイツェンベックねじれとスカラー $T$ は閉形式をもつ；
\item 動径ラピディティ~\eqref{eq:radial-rapidity} に対し、パラメーター空間不変量は
$J(\varphi)=\tfrac12\varphi(r)^{2}$ を満たし、場の恒等式
$T=\tfrac12 T_{\mathrm{bulk}}+ \nabla_{\mu}B^{\mu}$
は明示的に構成された動径カレント $B^{\mu}$ とともに成り立つ（静的・対角ゲージ切片）。
\end{enumerate}
\end{theorem}

\begin{conjecture}[一般モーター同一視]
チャート上の一般的な滑らかなモーター場 $M(x)=R(x)T(x)$ に対し、$(\Omega_{\mathrm{torsion}},\lambda)$ から構成される拡張診断 $J^{(5)}$ は、全発散までテレパラレルスカラー~\eqref{eq:teleparallel-T} と一致する：
\begin{equation}
J^{(5)} \;=\; \tfrac12\, T \;+\; \nabla_{\mu}B^{\mu}.
\label{eq:J5-T-conjecture}
\end{equation}
\end{conjecture}

この予想は $\lambda^{\mu}=0$ の静的対角シュワルツシルト切片上で上記定理に帰着する。任意モーターに対する証明には、ここで用いたチャート計算を超える $\Omega_{\mathrm{biv}}^{(5)}$ とヴァイツェンベックコントルジョンの完全な辞書が必要である。

\subsection{アインシュタイン--ヒルベルトとの関係}

積分恒等式 $\int T\,e\,d^{4}x=-\int R\sqrt{-g}\,d^{4}x+(\text{境界})$ およびそれに続くアインシュタイン方程式との変分的等価は TEGR において標準的である~\cite{maluf2013}。これらを外部入力として扱い、一般時空計算の再導出は本稿では行わない。$G(3,1,1)$ の退化計量全体に関するサンドイッチ、および $[\Omega_{\mathrm{torsion}},\Omega_{\mathrm{trans}}]\neq0$ のときの $M$ と $\exp(\Omega_{\mathrm{biv}}^{(5)})$ の同一視は、第5節と同様に開のままである。

\section{結論}

二重時空構造を、$\mathrm{Cl}(3,1)$ の双四元数ベクトル基底 $\{j, kI, kJ, kK\}$ に新しい零ベクトル $e_{4}$ を付加することにより射影幾何代数 $G(3,1,1)$ に埋め込んだ。拡張ねじれ不整合の生成子は、3つの双曲型生成子 $\{iI,iJ,iK\}$、3つの巡回型生成子 $\{I,J,K\}$、および4つの零バイベクトル $N_{\mu}=e_{4}\wedge e_{\mu}$（$\mu=0,1,2,3$）からなる10次元バイベクトル空間に整理される。これは $\mathfrak{iso}(3,1)\cong\mathfrak{so}(3,1)\ltimes\mathbb{R}^{3,1}$ の構造的候補である。

すべての $\mu,\nu$ に対する強い消去性 $N_{\mu}N_{\nu}=0$ により零部分の指数は1次の有限多項式となり、反転奇性 $\widetilde{N_{\mu}}=-N_{\mu}$ により積モーター $M:=RT$ のユニタリー性 $M\widetilde{M}=1$ が得られる。因子が可換でないとき、$\exp(\Omega_{\mathrm{biv}}^{(5)})$ との関係および退化計量全体のサンドイッチは開のままである。

パラメーター空間形式
\[
J^{(5)} \;=\; \tfrac{1}{16}B_{\text{Killing}}\!\bigl(\Omega_{\text{torsion}},\Omega_{\text{torsion}}\bigr) \;+\; \tfrac{1}{2}\eta_{\mu\nu}\lambda^{\mu}\lambda^{\nu}
\]
はねじれプラス並進係数上のスカラー診断として働く。双対は6次元ねじれバイベクトルセクターを閉じるが、各零生成子は次数4へ送られる。面に基づくブラケットと双対＝法線幾何は開プログラムとして残す（付録B）。第6節では固定チャート上の PGA--TEGR 対応を展開した：モーター誘導テトラッド、ヴァイツェンベックねじれ、シュワルツシルト対角 ansatz、および $J$ と $T$ のチャート水準の結びつきを述べ、一般の $J^{(5)}\leftrightarrow T$ 同一視は予想とし、TEGR／アインシュタイン--ヒルベルト変分等価は文献から取った。

\appendix

\section{10次元生成子の明示的対応}

拡張ねじれ不整合は3類に分解される10次元バイベクトル空間によって生成される。以下、親理論の双四元数ベクトル基底 $\{e_{0}=j,\ e_{1}=kI,\ e_{2}=kJ,\ e_{3}=kK\}$ と付加零ベクトル $e_{4}$（$e_{4}^{2}=0$、$\{e_{4},e_{\mu}\}=0$）を用いた $G(3,1,1)$ 設定における明示的代数代表を示す。

\subsection{双曲型生成子（平方 $+1$）}

次の3生成子はブースト様変換に対応する：

\begin{align*}
B_{1}^{+} &= iI = e_{0}\wedge e_{1}, \\
B_{2}^{+} &= iJ = e_{0}\wedge e_{2}, \\
B_{3}^{+} &= iK = e_{0}\wedge e_{3}.
\end{align*}

各々 $(B_{a}^{+})^{2} = +1$ を満たす。

\subsection{巡回型生成子（平方 $-1$）}

次の3生成子は双対時空の回転様変換に対応する：

\begin{align*}
B_{1}^{-} &= I = e_{3}\wedge e_{2}, \\
B_{2}^{-} &= J = e_{1}\wedge e_{3}, \\
B_{3}^{-} &= K = e_{2}\wedge e_{1}.
\end{align*}

各々 $(B_{a}^{-})^{2} = -1$ を満たす。

\subsection{零型生成子（平方 $0$）}

零生成子は付加零ベクトル $e_{4}$ と4つの $\mathrm{Cl}(3,1)$ ベクトル生成子から構成される4つのバイベクトルである：
\begin{align*}
N_{0} &= e_{4}\wedge e_{0} = e_{4}\,j,         &\quad& \text{（時間並進生成子）}\\
N_{1} &= e_{4}\wedge e_{1} = e_{4}\,(kI),      &\quad& \text{（$x$ 並進生成子）}\\
N_{2} &= e_{4}\wedge e_{2} = e_{4}\,(kJ),      &\quad& \text{（$y$ 並進生成子）}\\
N_{3} &= e_{4}\wedge e_{3} = e_{4}\,(kK),      &\quad& \text{（$z$ 並進生成子）}
\end{align*}
$e_{4}^{2}=0$ および $\{e_{4},e_{\mu}\}=0$ を用いると、\emph{任意} の $\mu,\nu\in\{0,1,2,3\}$ に対して直接
\[
N_{\mu}\,N_{\nu} \;=\; (e_{4}e_{\mu})(e_{4}e_{\nu}) \;=\; -\,e_{4}^{2}\,e_{\mu}e_{\nu} \;=\; 0.
\]
したがって各 $N_{\mu}$ の平方は零（$\mu=\nu$）であり、任意の組の積も消える；4つの $N_{\mu}$ はポアンカレ代数 $\mathfrak{iso}(3,1)$ の並進生成子 $P_{\mu}$ を実現する可換な並進イデアルを形成する。反転は $\widetilde{N_{\mu}}=-N_{\mu}$ であり、モーターのユニタリー性（第5節）に不可欠である。

\paragraph{通常時空の光様ベクトルとの関係。}
親理論の光様ベクトルとの接続は直接的である。通常時空の一般的光様ベクトルは
\[
L \;=\; ct\,j + x\,kI + y\,kJ + z\,kK \;=\; ct\,e_{0} + x\,e_{1} + y\,e_{2} + z\,e_{3},
\]
ミンコフスキー符号の零条件
\[
L^{2} \;=\; -(ct)^{2} + x^{2} + y^{2} + z^{2} \;=\; 0,
\]
（親論文の符号 $(-,+,+,+)$）を満たす。この方向 $L$ \emph{に沿った}並進は零バイベクトル
\[
T_{L} \;:=\; e_{4}\wedge L \;=\; ct\,N_{0} + x\,N_{1} + y\,N_{2} + z\,N_{3},
\]
によって生成される。これは4つの零バイベクトルの線形結合であり、$N_{\mu}N_{\nu}$ の強い消去により自身の平方も零である。並進方向の光様性は生成子そのものへの代数的制約ではなく、4つの並進パラメーター $\lambda^{\mu}=(ct,x,y,z)$ 上の錐条件
\[
\eta_{\mu\nu}\,\lambda^{\mu}\lambda^{\nu} \;=\; -(ct)^{2}+x^{2}+y^{2}+z^{2} \;=\; 0
\]
として表される。4つの $N_{\mu}$ は線形独立であり、並進セクターは真に4次元のままである。

\subsection{まとめ}

10生成子全体とその代数性質を表~\ref{tab:generators} にまとめる。

\begin{table}[htbp]
\centering
\caption{10次元生成子構造}
\label{tab:generators}
\begin{tabular}{llcl}
\hline
類           & 生成子                                                                & 個数 & 平方 \\
\hline
双曲型       & \(B_{1}^{+}=iI,\ B_{2}^{+}=iJ,\ B_{3}^{+}=iK\)                        & 3     & \(+1\) \\
巡回型       & \(B_{1}^{-}=I,\ \ B_{2}^{-}=J,\ \ B_{3}^{-}=K\)                       & 3     & \(-1\) \\
零型         & \(N_{0}=e_{4}j,\ N_{1}=e_{4}(kI),\ N_{2}=e_{4}(kJ),\ N_{3}=e_{4}(kK)\) & 4     & \(\phantom{+}0\) \\
\hline
\end{tabular}
\end{table}

この10次元基底は拡張ねじれ不整合バイベクトル $\Omega_{\text{biv}}^{(5)}$ の生成子空間を形成する。双対 $X\mapsto Xi$ は6つの双曲・巡回生成子をバイベクトル次数の内部で閉じ（かつ $e_{4}$ は双四元数擬スカラー $i$ と可換である）、一方各零生成子 $N_{\mu}$ は次数4へ送られる。並進は $e_{4}$ と通常時空ベクトル生成子から内在的に生成され、光様並進は4次元並進パラメーター空間の錐 $\eta_{\mu\nu}\lambda^{\mu}\lambda^{\nu}=0$ に対応する。

\section{面に基づくブラケット積と双対時空対応の詳細}

本付録では、第4節の面に基づくブラケット積と双対時空の法線性質との対応を補足する明示計算を与える。

\subsection{ブラケット積の定義と次数内容}

$G(3,1,1)$ 内の2つの3ブレード $A,B$ に対し、面に基づくブラケット積を
\[
A \mathbin{\underline{\vee}} B \;=\; \tfrac12 (AB - BA), \qquad
A \mathbin{\overline{\vee}} B \;=\; \tfrac12 (AB + BA).
\]
と定義する。第4節（脚注）で説明したとおり、これらの記号は \eqref{eq:bracket-def} の意味でのブラケット積であり、反スカラー単位による標準PGAの「反積」とは一致しない。

5ベクトル代数 $G(3,1,1)$ では、2つの3ブレードの幾何積の次数内容は
\[
A_{3}B_{3} \;=\; \langle A_{3}B_{3}\rangle_{0} \;+\; \langle A_{3}B_{3}\rangle_{2} \;+\; \langle A_{3}B_{3}\rangle_{4}.
\]
入替 $A\leftrightarrow B$ の下で各次数成分は確定した対称性をもつと期待される：
\begin{equation}
\langle A_{3}B_{3}\rangle_{0} = +\langle B_{3}A_{3}\rangle_{0},
\quad
\langle A_{3}B_{3}\rangle_{2} = -\langle B_{3}A_{3}\rangle_{2},
\quad
\langle A_{3}B_{3}\rangle_{4} = +\langle B_{3}A_{3}\rangle_{4}.
\label{eq:grade-symmetry}
\end{equation}

\paragraph{明示例による検証。}
$A=e_{1}e_{2}e_{3}$、$B=e_{1}e_{2}e_{4}$ を $G(3,1,1)$ でとる。$e_{1}^{2}=e_{2}^{2}=e_{3}^{2}=+1$、$e_{4}^{2}=0$、生成子はすべて反可換。$e_{3}$ を $e_{1}e_{2}$ の向こうへ反可換させ $e_{1}e_{2}e_{1}e_{2}=-1$ を用いると、直接計算で
\[
AB \;=\; e_{1}e_{2}e_{3}\,e_{1}e_{2}e_{4} \;=\; -\,e_{3}e_{4},
\]
これは純次数2である。同様に
\[
BA \;=\; e_{1}e_{2}e_{4}\,e_{1}e_{2}e_{3} \;=\; -\,e_{4}e_{3} \;=\; +\,e_{3}e_{4}.
\]
したがって
\[
A\mathbin{\underline{\vee}}B \;=\; \tfrac{1}{2}(AB-BA) \;=\; -\,e_{3}e_{4} \;\in\; \text{次数2},
\quad
A\mathbin{\overline{\vee}}B \;=\; \tfrac{1}{2}(AB+BA) \;=\; 0,
\]
となり \eqref{eq:grade-symmetry} と一致する。

\paragraph{非零次数4の場合の検証。}
$A=e_{1}e_{2}e_{3}$、$B=A=e_{1}e_{2}e_{3}$ とする。$AB=A^{2}=(e_{1}e_{2}e_{3})^{2}=-1$ は次数0；$BA=A^{2}=-1$；よって $A\mathbin{\overline{\vee}}A=-1$（スカラー）かつ $A\mathbin{\underline{\vee}}A=0$。次数4の例として $A=e_{0}e_{1}e_{2}$、$B=e_{0}e_{3}e_{4}$ をとる。$AB=(e_{0}e_{1}e_{2})(e_{0}e_{3}e_{4})$ の次数4成分は非零の $e_{1}e_{2}e_{3}e_{4}$（生成子の順序に依存する関連4ブレード）である。これらの計算は考察した例における次数パターンを支持する；任意の3ブレードに対する一般定理は開のままである。

\subsection{双対時空法線との対応}

$P$ を通常時空ベクトル生成子から構成した3ブレード、$D$ を双対時空セクターの元とする。代表的な構成では
\begin{equation}
P \mathbin{\underline{\vee}} D \;=\; \kappa\,(e_{4}\wedge n),
\label{eq:appB-normal}
\end{equation}
が成り立つ。ここで $n$ は $P$ が表す面の（通常時空の意味での）法線ベクトルである。
また $e_{4}\wedge n$ は $\{N_{0},N_{1},N_{2},N_{3}\}$ が張る4次元零セクター内の零バイベクトルであり、
$\kappa$ は $P$ と $D$ の正規化から定まる実係数である。両辺は次数2であり \eqref{eq:grade-symmetry} と整合する。

\paragraph{具体例。}
$P = e_{4}\wedge e_{1}\wedge e_{2}$（$e_{4}$ を1因子もち $e_{1},e_{2}$ が張る2平面を含む3ブレード）および $D=jI$（親理論の双対時空セクターの元）をとる。$jI=-e_{0}e_{3}e_{2}$（親理論の双対基底の次数3同定）および $\{e_{4},e_{\mu}\}=0$ より
\[
PD \;=\; (e_{4}e_{1}e_{2})\,(-e_{0}e_{3}e_{2}) \;=\; -\,e_{4}e_{1}e_{2}e_{0}e_{3}e_{2}.
\]
$e_{2}$ を自身の向こうへ反可換させ因子を消去（$e_{2}^{2}=+1$）し、並べ替えると $PD-DP$ は非零の $e_{4}\wedge e_{3}$ の倍数、すなわち零セクター内の $N_{3}=e_{4}\wedge e_{3}$ に比例する元に帰着する。方向 $e_{3}$ は（空間部分において）2平面 $e_{1}\wedge e_{2}$ の単位法線であるから、$\kappa\cdot(e_{4}\wedge n) = \kappa\,N_{3}$、$n=e_{3}$ となり、この代表で \eqref{eq:appB-normal} が確認される。

\eqref{eq:appB-normal} の右辺は「$e_{4}\wedge(\text{ベクトル})$」形のバイベクトル、すなわち零並進イデアルの元である。双対時空の元 $D$ はしたがって法線方向の担体として作用し、ブラケット $P\mathbin{\underline{\vee}}D$ はその法線に沿った並進の生成子を生む。

\subsection{並進不整合の出現}

3ブレード $P$ を零セクターからの寄与で拡大すると、
\[
P \;\to\; P \;+\; \sum_{\mu=0}^{3}\lambda^{\mu}\,P^{(\mu)}_{\text{null}},
\qquad
P^{(\mu)}_{\text{null}} \;:=\; N_{\mu}\wedge n,
\]
反交換子ブラケット $P\mathbin{\overline{\vee}}D$ に並進パラメーターに線形な追加項が現れる：
\[
P \mathbin{\overline{\vee}} D
\;=\; \bigl(P\mathbin{\overline{\vee}}D\bigr)_{\text{torsional}}
\;+\; \sum_{\mu=0}^{3} c_{\mu}\,\lambda^{\mu}\,(e_{4}\wedge n\wedge e_{\mu}),
\]
次数4の4ブレード $e_{4}\wedge n\wedge e_{\mu}$ が並進不整合の幾何的担体である。係数 $c_{\mu}$ は面と双対法線の相対的向きで定まる。すべて $\lambda^{\mu}=0$ のとき並進項は恒等的に消え純ねじれの場合に戻る。特に $e_{4}^{2}=0$ により消える $N_{\mu}\cdot e_{4}$ 型の縮約はここには現れない；正しい幾何対象は次数4の4ブレード $e_{4}\wedge n\wedge e_{\mu}$ である。

\subsection{双対の下での整合性}

双対写像 $X\mapsto Xi$ は $\mathrm{Cl}(3,1)$ 部分代数上で標準的に作用し、$e_{4}i=ie_{4}$ により $e_{4}$ 因子を通り抜ける。零バイベクトル $N_{\mu}$ の双対化は次数4の元を与える（第2節）ため、零セクター上の完全に双対共変なブラケット計算は開のままである。

これらの計算はねじれ・並進不整合の幾何的読みを例示し、第4節の一般プログラムは後続の課題とする。

\begin{thebibliography}{99}

\bibitem{Gunn2017}
C.~Gunn,
\textit{Geometric algebras for Euclidean geometry},
Advances in Applied Clifford Algebras \textbf{27} (1), 185--208 (2017).

\bibitem{Lengyel2024}
E.~Lengyel,
\textit{Projective Geometric Algebra Illuminated},
Terathon Software LLC (2024).

\bibitem{DeKeninckDorst2019}
S.~De Keninck and L.~Dorst,
\textit{Projective geometric algebra: A new framework for doing Euclidean geometry},
arXiv:1901.05873 (2019).

\bibitem{GunnDeKeninck2020}
C.~Gunn and S.~De Keninck,
\textit{Geometric algebra and computer graphics},
in: SIGGRAPH 2019 Courses, ACM, New York (2019).
See also: \texttt{https://bivector.net/}.

\bibitem{DorstFontijneMann2007}
L.~Dorst, D.~Fontijne, and S.~Mann,
\textit{Geometric Algebra for Computer Science: An Object-Oriented Approach to Geometry},
Morgan Kaufmann, San Francisco (2007).

\bibitem{DoranLasenby2003}
C.~Doran and A.~Lasenby,
\textit{Geometric Algebra for Physicists},
Cambridge University Press, Cambridge (2003).

\bibitem{Hestenes2015}
D.~Hestenes,
\textit{Space-Time Algebra}, 2nd ed.,
Birkh\"{a}user, Basel (2015).

\bibitem{HestenesSobczyk1984}
D.~Hestenes and G.~Sobczyk,
\textit{Clifford Algebra to Geometric Calculus: A Unified Language for Mathematics and Physics},
Reidel, Dordrecht (1984).

\bibitem{Lounesto2001}
P.~Lounesto,
\textit{Clifford Algebras and Spinors}, 2nd ed.,
London Mathematical Society Lecture Note Series \textbf{286},
Cambridge University Press, Cambridge (2001).

\bibitem{ConwaySmith2003}
J.~H.~Conway and D.~A.~Smith,
\textit{On Quaternions and Octonions: Their Geometry, Arithmetic, and Symmetry},
A K Peters, Natick, MA (2003).

\bibitem{Selig2005}
J.~M.~Selig,
\textit{Geometric Fundamentals of Robotics}, 2nd ed.,
Monographs in Computer Science, Springer, New York (2005).

\bibitem{Weinberg1995}
S.~Weinberg,
\textit{The Quantum Theory of Fields, Volume I: Foundations},
Cambridge University Press, Cambridge (1995).

\bibitem{Tung1985}
W.-K.~Tung,
\textit{Group Theory in Physics},
World Scientific, Singapore (1985).

\bibitem{FultonHarris1991}
W.~Fulton and J.~Harris,
\textit{Representation Theory: A First Course},
Graduate Texts in Mathematics \textbf{129},
Springer, New York (1991).

\bibitem{Humphreys1972}
J.~E.~Humphreys,
\textit{Introduction to Lie Algebras and Representation Theory},
Graduate Texts in Mathematics \textbf{9},
Springer, New York (1972).

\bibitem{maluf2013}
J.~W.~Maluf,
\textit{The teleparallel equivalent of general relativity},
Annalen der Physik \textbf{525} (5--6), 339--357 (2013).

\end{thebibliography}

\end{document}
